\section{Related Work}

ARCH and GARCH models provide the classical foundation for conditional volatility forecasting. Engle introduced ARCH to model time-varying conditional variance, and Bollerslev generalized the recursion to include lagged conditional variance \cite{engle1982arch,bollerslev1986garch}. The empirical appeal of GARCH(1,1) is not only parsimony. Large-scale forecast comparisons have shown that it can be difficult to beat in practical volatility forecasting applications \cite{hansen2005forecast}.

Mean dynamics are often modeled before conditional variance. The Box-Jenkins ARIMA framework captures serial dependence in univariate time series \cite{box2015time}. In financial returns, the conditional mean is usually small relative to the conditional variance, but an ARMA or ARIMA mean model can still be paired with GARCH to test whether residual variance forecasts improve after filtering predictable mean components.

GARCH-family extensions address stylized features that a symmetric GARCH(1,1) recursion may miss. EGARCH models allow asymmetric volatility responses and avoid some positivity restrictions \cite{nelson1991egarch}. GJR-type models capture sign-dependent shock effects \cite{glosten1993gjr}. APARCH and FIGARCH specifications allow richer power transformations and long-memory volatility dynamics \cite{ding1993aparch,baillie1996figarch}. These extensions motivate the broad specification search evaluated in this paper.

Neural sequence models offer a different forecasting premise. LSTM networks were designed to learn long-range sequential dependencies through gated recurrent states \cite{hochreiter1997lstm}. In volatility forecasting, LSTMs can use lagged returns, volume, rolling volatility, and econometric forecasts as supervised features. Hybrid econometric-neural designs are therefore useful empirical candidates, but their performance must be judged against strong volatility-specific baselines rather than against weak naive alternatives.

Forecast evaluation is central because latent volatility is not directly observed. Squared daily returns are available and simple, but they are noisy proxies for conditional variance. QLIKE is widely used because it remains attractive when forecast comparisons rely on imperfect volatility proxies and because it evaluates forecasts on the variance scale \cite{patton2011volatility}. High-frequency realized volatility can provide richer targets when intraday data are available \cite{andersen2003realized}, but daily OHLC estimators are also useful. Parkinson, Garman-Klass, Rogers-Satchell, and Yang-Zhang estimators exploit high, low, open, and close information to obtain alternative volatility proxies \cite{parkinson1980extreme,garman1980estimation,rogers1991estimating,yang2000drift}.

Statistical forecast comparisons require care when many models are examined. Diebold-Mariano tests compare predictive accuracy under a specified loss difference process \cite{diebold1995predictive}. Holm correction controls family-wise error rates across multiple tests \cite{holm1979sequential}. The model confidence set (MCS) framework provides a complementary way to identify a set of models not statistically eliminated under a loss criterion \cite{hansen2011mcs}. VaR backtesting adds a risk-management view through unconditional and conditional coverage tests \cite{kupiec1995var,christoffersen1998interval}.
